\chapter{Як машина вчиться рухатися}
% full version: chapters/15_motorlearning.tex

Ми бачили два способи вчитися. Без учителя --- «разом, отже, пов'язані» --- машина запам'ятовує те, що часто трапляється. З дофаміном вона вчиться на сюрпризах: вийшло краще чи гірше, ніж очікувалося. Але коли ви схибили, кидаючи м'яч у кошик, ви знаєте не лише, що вийшло погано. Ви знаєте, \emph{наскільки} і \emph{в який бік} схибили. Це третій учитель --- помилка з напрямком, особиста для кожного виходу.

\section*{Провід помилки}

У мозку є окрема схема для навчання рухів, і влаштована вона дуже незвично. Величезна кількість дрібних \nt{GLUT}-нейронів, десятки мільярдів, розкладає вхідні сигнали в дуже довгий список ознак. Вихідні нейрони цієї схеми --- гальмівні, \nt{GABA}, їх десятки мільйонів, і кожен слухає сотню тисяч входів. І до кожного вихідного нейрона тягнеться рівно один особливий провід --- провід помилки. Якщо помилка прийшла, поки вхід був активний, цей вхід слабшає. Це класичне правило «зменшуй помилку», яким інженери навчають адаптивні фільтри.

\pencilfig{15}{\pencilart{15}}{Промах повідомляє не лише «погано», а й «наскільки і в який бік».}

Такий учитель швидший за дофаміновий: у кожного виходу своя помилка, і не треба вивуджувати свою частку із загального сигналу на всіх. Щоправда, помилка приходить із запізненням, і з'єднанню знову потрібен «стікер», що пам'ятає, що воно брало участь. Досліди показують, що тривалість цього стікера підібрано під затримку помилки.

\section*{Внутрішня модель}

Чого навчається така схема? Моделі власного тіла: який результат дасть команда. З нею не треба чекати запізнілих датчиків --- результат можна передбачити заздалегідь. За однією з гіпотез, саме тому ви не можете полоскотати самі себе: мозок заздалегідь знає, що буде, і віднімає передбачене.

\section*{Швидкий і повільний учень}

Вдягніть окуляри, що зсувають картинку вбік, і кидайте дротики. Перші кидки йдуть убік, але за десятки спроб ви пристосуєтеся. Зніміть окуляри --- і знову промахи, тепер в інший бік. Цікаве починається далі. Досліди добре описуються двома учнями одразу: швидкий швидко вчиться і швидко забуває, повільний учиться повільно, але пам'ятає довго. Якщо після довгого пристосування коротко перевчитися назад, загальна поправка повернеться до нуля --- але старе вміння потім саме частково спливе: швидкий учень забув, а повільний пам'ятає. Модель з одним учнем так не вміє, а люди --- вміють.

\pencilfig{115}{%
% figures/tworate_*.dat from models/motor_learning.py: adapt 200 throws, reverse until the sum is back to zero, then no errors at all
\begin{scope}[x=0.026cm, y=1.45cm]
\fill[pcGray, opacity=0.08] (220,-1.1) rectangle (226,1.1);
\draw[pen] (0,0) -- (340,0);
\draw[pcGray, line width=0.4pt] plot file {figures/tworate_p.dat};
\draw[penlite=pcAmber] plot file {figures/tworate_fast.dat};
\draw[penlite=pcBlue] plot file {figures/tworate_slow.dat};
\draw[pcGray!40!black, line width=1.1pt] plot file {figures/tworate_net.dat};
\draw[pcGray, densely dashed, line width=0.4pt] (226,-1.1) -- (226,1.1);
\end{scope}
\node[lbl, anchor=south] at (3.1,1.47) {окуляри вдягнуто};
\node[lbl, anchor=north west, align=left] at (6.3,-0.55) {коротко\\перевчилися};
\node[lbl, anchor=south west] at (6.0,1.47) {помилок немає};
\node[lbl, text=pcBlue, anchor=north] at (4.4,0.8) {повільний};
\node[lbl, text=pcAmber!80!black, anchor=south west] at (1.15,0.5) {швидкий};
\node[lbl, text=pcGray!40!black, anchor=south] at (7.7,0.95) {сума спливла};
}{Поправка двох учнів у моделі. Після короткого перевчання сума на нулі, але повільний пам'ятає старе, і воно саме спливає.}

Навіщо два учні? За гіпотезою, це той самий рецепт, що в штурмана з дев'ятого розділу: світ змінюється і швидко, і повільно, і розумно стежити за обома.

\fullversion{15}
